\begin{tikzpicture}[x=1cm,y=1cm,font=\footnotesize]
  \path[use as bounding box] (0,0) rectangle (17.2,5.15);

  % Panel frames and titles
  \draw[black!25,line width=0.4pt] (0.05,0.10) rectangle (8.25,5.00);
  \draw[black!25,line width=0.4pt] (8.70,0.10) rectangle (17.15,5.00);
  \node[anchor=west,font=\bfseries] at (0.25,4.72)
    {(a) Transit geometry};
  \node[anchor=west,font=\bfseries] at (8.90,4.72)
    {(b) Wavelength-dependent transit depth};

  % Star
  \shade[inner color=yellow!15,outer color=transitstar]
    (1.60,2.60) circle (1.23);
  \draw[transitstar!70!black,line width=0.5pt]
    (1.60,2.60) circle (1.23);
  \node[anchor=north] at (1.60,1.20) {Host star};

  % Wavelength-labelled stellar rays before and after the atmosphere
  \draw[transitstar!80!black,line width=1.15pt]
    (2.72,3.10) -- (4.10,3.10);
  \draw[transitabsorb,line width=0.55pt,->]
    (5.32,3.10) -- (7.72,3.10);
  \node[anchor=south,text=transitabsorb] at (3.42,3.13)
    {$\lambda_{\rm high}$};

  \draw[transitstar!80!black,line width=1.15pt]
    (2.72,2.02) -- (4.16,2.02);
  \draw[transitatmos,line width=0.95pt,->]
    (5.27,2.02) -- (7.72,2.02);
  \node[anchor=north,text=transitatmos] at (3.43,1.98)
    {$\lambda_{\rm low}$};

  % Opaque planet and wavelength-dependent atmospheric radii
  \fill[transitatmos!12] (4.72,2.60) circle (0.78);
  \draw[transitabsorb,dashed,line width=1.05pt]
    (4.72,2.60) circle (0.78);
  \draw[transitatmos,line width=1.05pt]
    (4.72,2.60) circle (0.61);
  \fill[transitplanet] (4.72,2.60) circle (0.45);
  \draw[black,line width=0.45pt] (4.72,2.60) circle (0.45);

  \draw[transitabsorb,dashed,line width=1.05pt]
    (5.30,4.08) -- (5.86,4.08);
  \node[anchor=west,text=transitabsorb] at (5.96,4.08)
    {larger $R_{\rm p}(\lambda_{\rm high})$};
  \draw[transitatmos,line width=1.05pt]
    (5.30,3.72) -- (5.86,3.72);
  \node[anchor=west,text=transitatmos] at (5.96,3.72)
    {smaller $R_{\rm p}(\lambda_{\rm low})$};

  \draw[black!60,->,line width=0.55pt]
    (5.82,2.59) -- (5.32,2.67);
  \node[anchor=west,align=left] at (5.90,2.58)
    {atmospheric\\annulus};

  \node[anchor=north,align=center] at (4.72,1.58)
    {Planet and\\atmospheric annulus};
  \node[anchor=north,align=center] at (7.35,1.88)
    {toward the\\observer};

  % Spectrum axes
  \draw[->,line width=0.65pt] (9.45,0.83) -- (16.78,0.83);
  \draw[->,line width=0.65pt] (9.45,0.83) -- (9.45,4.38);
  \node[anchor=north] at (13.10,0.66) {Wavelength, $\lambda$};
  \node[rotate=90,anchor=south] at (8.92,2.63)
    {Transit depth, $\delta(\lambda)$};

  % Illustrative transmission spectrum
  \draw[black!15,densely dotted] (9.45,1.43) -- (16.64,1.43);
  \draw[black!20,densely dotted] (9.45,2.75) -- (16.64,2.75);
  \draw[transitplanet,line width=1.35pt]
    plot[smooth,tension=0.72] coordinates {
      (9.52,1.47) (9.88,1.50) (10.22,1.62) (10.56,2.15)
      (10.90,1.70) (11.28,1.52) (11.65,1.58) (12.03,1.86)
      (12.38,3.05) (12.73,1.91) (13.08,1.56) (13.48,1.49)
      (13.88,1.68) (14.27,2.19) (14.65,1.73) (15.03,1.62)
      (15.39,2.00) (15.76,3.48) (16.10,2.03) (16.50,1.57)
    };

  % Link spectral extrema to the two atmospheric radii in panel (a)
  \fill[transitabsorb] (12.38,3.05) circle (0.065);
  \draw[transitabsorb,densely dashed] (12.38,0.84) -- (12.38,3.05);
  \node[anchor=south,text=transitabsorb,align=center]
    at (12.38,3.18) {higher opacity\\deeper transit};

  \fill[transitatmos] (13.48,1.49) circle (0.065);
  \draw[transitatmos,densely dashed] (13.48,0.84) -- (13.48,1.49);
  \node[anchor=west,text=transitatmos,align=left]
    at (13.63,1.44) {lower opacity\\shallower transit};

  \node[anchor=east,align=right,font=\scriptsize] at (16.68,4.20)
    {$\mathrm{opacity}\!\uparrow\;\Rightarrow\;
      R_{\rm p}(\lambda)\!\uparrow\;\Rightarrow\;
      \delta(\lambda)\!\uparrow$};
\end{tikzpicture}
